With the increasing number of undergraduate theses and academic papers, the standardization and efficiency of paper typesetting have attracted growing attention. Traditional word processors are easy to use, but they still require much manual adjustment in complex formulas, figure and table numbering, cross-references, reference management, and template adaptation. Although LaTeX is powerful in academic typesetting, its syntax, environment structures, compilation configuration, and error logs are not beginner-friendly. To address problems such as high code memorization cost, complex template usage, and difficult error locating, this thesis designs and implements an intelligent academic paper typesetting system based on JavaWeb and LaTeX, providing online, component-based, and intelligent support for paper writing and typesetting.

The system adopts a front-end and back-end separated architecture. The front end is developed with Vue 3 and Vite, while the back end uses Spring Boot, MyBatis, and MySQL for business logic processing and data persistence. The system provides functions such as user registration and login, paper project management, online LaTeX editing, template management, paper compilation, PDF preview, file export, and AI assistance. Users can create paper projects in the browser, select and import templates, edit papers online, and generate PDF documents through different compilation engines.

The core design of the system is a component library-based drag-and-drop insertion mechanism for paper structures. Common LaTeX structures, including chapters, images, tables, formulas, and symbols, are encapsulated as visual components. Users do not need to memorize complete LaTeX commands or environment syntax. They only need to drag required components into the editing area, and the system automatically generates corresponding code frameworks. This approach lowers the threshold for using LaTeX, reduces command misspellings, missing brackets, and unclosed environments, and improves paper editing efficiency.